\section{Introduction}
\label{sec:intro}

High-frequency bus systems face a persistent coordination challenge.
Timetables are planned offline to balance fleet utilisation against passenger demand, yet the headways they prescribe degrade within minutes once buses encounter stochastic traffic and boarding delays.
Real-time holding control can restore regularity, but only within the headway structure it inherits.
When the timetable itself is poorly matched to demand, dispatching too few buses in the peak or too many in the off-peak, holding alone may not recover acceptable service.
Planning and control are thus coupled: the quality of one layer bounds the achievable performance of the other.

Despite this coupling, existing work treats the two problems separately.
The timetable optimisation literature formulates headway or frequency design as mixed-integer programmes solved offline \citep{gkiotsalitis2021bus}, while the holding control literature trains reinforcement learning (RL) agents that take the timetable as given \citep{wang2020real,zhang2026single}.
Joint optimisation would seem a natural application of hierarchical reinforcement learning (HRL), yet three gaps prevent a straightforward application.

\paragraph{Asynchrony between planning and control.}
HRL frameworks such as the options framework \citep{sutton1999between} and goal-conditioned methods such as HIRO \citep{nachum2018data} assume that the upper level sets subgoals or options that can be evaluated through a lower-level execution segment.
In bus operations, the upper policy acts at dispatch events (once per bus departure), while the lower policy acts at every station arrival across all active buses.
The two levels observe different state spaces: the upper level sees route-level demand forecasts and fleet status, whereas the lower level sees per-bus headway deviations and passenger loads.
They also interact through delayed feedback spanning an entire trip lifecycle.
More importantly, several dispatch goals are active concurrently: one upper decision may set the planned launch time of a future trip while buses dispatched under earlier goals continue to move and hold at stations.
Thus the issue is not only the absence of a shared clock, but also the need to optimise and credit overlapping timetable edits.

\paragraph{Cross-advantage under asynchrony.}
Pan et al.~\citep{pan2024roboduet} recently introduced cross-advantage augmentation for bi-level robot locomotion, where an upper style policy and a lower tracking policy execute synchronously on the same state.
Their augmented advantage $\AUaug = \AU + \beta \AL$ directly sums the two advantages at each time step.
This mechanism is not directly applicable when the two levels operate on different timescales with different state and action spaces: no single time step jointly defines $\AU$ and $\AL$.
Generalising cross-level coupling to this asynchronous, multi-timescale setting requires propagating execution outcomes over the variable-length trip lifecycle induced by each dispatch.
It also requires exposing lower-level holding statistics to the upper policy through state augmentation rather than reward backflow.

\paragraph{Fleet and regularity constraints.}
Bus operations also impose feasibility-style constraints, including a fleet-size budget $\Nfleet$ and bounded headway deviation.
Fixed-coefficient reward shaping may not accommodate these constraints reliably.
Constrained RL methods such as reward constrained policy optimisation (RCPO) \citep{tessler2019reward} address single-level constraint satisfaction but do not account for constraints that span a bi-level structure, where the upper level controls resource allocation and the lower level controls service regularity.

\medskip

We propose TransitDuet, a bi-level reinforcement learning framework designed to address these three gaps.
The upper policy $\piU$ selects a discrete planned-headway adjustment for the next same-direction trip; the simulator converts the resulting planned headway into a planned launch time, so the upper directly modifies the rolling timetable. The lower policy $\piL$ is a soft actor-critic (SAC) with Lagrangian penalties \citep{haarnoja2018soft} that selects holding times at each station arrival, with parameters shared across buses.
Two auxiliary upper-only signals---a holding-feedback state channel and a per-dispatch hindsight credit---provide cross-level credit assignment without injecting upper signals into the lower transitions.
Fleet-size feasibility is maintained through $\thetavec$-weighted online gradient descent on a measurement vector, inspired by the ApproPO framework \citep{miryoosefi2019reinforcement}, while headway uniformity enters the lower-level objective through Lagrangian relaxation.

\paragraph{Contributions.}
The specific contributions of this paper are:
\begin{enumerate}
    \item Dynamic timetable--holding formulation. We cast rolling bus timetable optimisation and holding control as a bi-level Markov decision process with heterogeneous state-action spaces and asynchronous transitions.
    \item Discrete timetable and holding policies. The upper policy chooses from a finite dispatch-headway adjustment grid, and the lower policy chooses from a finite holding-time grid augmented with adjacent planned-dispatch context. This discretisation matches operational command granularity and isolates the value of dynamic timetable decisions against fixed-headway baselines.
    \item Adaptive constraint handling via $\thetavec$-weighted online gradient descent (OGD) and Lagrangian relaxation. At the upper level, measurement-driven OGD adjusts the reward weight on fleet overshoot. At the lower level, a dual variable penalises headway deviation. OGD removes the need to hand-tune an overshoot coefficient but provides no worst-case feasibility guarantee.
    \item Empirical evaluation on a calibrated corridor with ten random seeds, reporting both service-quality metrics (waiting time, headway CV, total passenger time) and timetable-execution diagnostics (planned-headway variability, dispatch lateness).
\end{enumerate}

\paragraph{Results preview.}
We evaluate on a high-frequency bus corridor with 22 stations, 264 daily trips, demand stochasticity $\sigma_{\mathrm{demand}} = 0.15$, and elastic fleet $\Nfleet \in [8, 16]$ sampled per episode. All learned methods use the same simulator and random-seed protocol, with the target-headway-only model reported as an ablation rather than as the main TransitDuet configuration.
